\section{What the final theorem does not use}\label{sec:offpath}

Nearly every file of the archive imports nearly every other, so the import graph says nothing
about what depends on what. The statements of this section are instead computed from the
compiled archive: the \emph{cone} of \lean{riemannHypothesis} is the set of all declarations
reachable from it through the type and the value of every declaration reached, the link
between an inductive type and its constructors, and the structure of a projection
(\cref{sec:cone}). It has \conesize{} declarations: \conearchive{} of the archive and
\ifsecondrun
\conemathlib{} of Mathlib and \coneother{} of Lean's library and the other packages.
\else
\conerest{} of Mathlib, of the packages it depends on and of Lean's own library.
\fi
The other four theorems of \lean{Main.lean} are derived from the same statement for all
characters, by short arguments of their own: the cone of all five has \conefive{}
declarations, \conefiveextra{} more.

\begin{table}[ht]
\centering
\begin{tabular}{@{}lrrr@{}}
\toprule
 & in the archive & in the cone & share\\
\midrule
Theorems & \projtheorems & \coneprojtheorems & \theoremshare\\
Definitions & \projdefs & \coneprojdefs & \defshare\\
Modules that declare something & \projmodules & \coneprojmodules & \moduleshare\\
\bottomrule
\end{tabular}
\caption{The archive and the cone of \texttt{riemannHypothesis}, counted by distinct names. A
module is counted in the cone if it declares at least one declaration of the cone. The
definitions include \projunsaferec{} compiled forms of recursive definitions, none of them in
the cone.}\label{tab:cone}
\end{table}

So about a third of the archive's theorems, \unusedtheorems{} of them, are not used by the
final theorem. They are compiled and checked by the kernel like the rest
(\cref{sec:verification}); they are simply not on the path. Of the archive's \filecount{}
files, \unusedmodules{} declare nothing that is in the cone%
\ifsecondrun
. By top-level directory they number
\modsDevelopmentUnused{} of \modsDevelopmentAll{} in \lean{Development},
\modsFormalismUnused{} of \modsFormalismAll{} in \lean{Formalism}, \modsGroundTheoryUnused{}
of \modsGroundTheoryAll{} in \lean{GroundTheory}, \modsReciprocityTheoryUnused{} of
\modsReciprocityTheoryAll{} in \lean{ReciprocityTheory}, \modsCompanionTheoryUnused{} of
\modsCompanionTheoryAll{} in \lean{CompanionTheory}, \modsPositivityTheoryUnused{} of
\modsPositivityTheoryAll{} in \lean{PositivityTheory}, and \modsOtherUnused{} in the
\modsOtherRoots{} small directories that remain; these totals count every file,
including the \emptymodules{} that declare nothing at all.
\else
, the \emptymodules{} that declare nothing at all among them.
\fi

What is unused includes earlier and parallel forms of steps that are on the path. Each of the
following was confirmed outside the cone by its full name.
\begin{itemize}
\item \emph{Conditional reductions.} Several statements of the form ``if such data exist at
  every scale, then the Riemann Hypothesis holds'':
  \begin{center}\small
  \begin{tabular}{@{}ll@{}}
  \lean{RiemannHypothesisOfLowerBounds} & \lean{RiemannHypothesisOfHalfLinks}\\
  \lean{WeightedSieveImpliesRH} & \lean{WeightedSieveImpliesRHThreeParameter}\\
  \lean{GrowthBoundGivesRiemannHypothesis} & \lean{RiemannHypothesisOfBalanceAndGrowthData}
  \end{tabular}
  \end{center}
  One more of this kind, \lean{ZetaQuarterPowerCorollary}, takes the number of lifts
  $m\approx x^{1/4}$ and concludes only that $\zeta(s)\neq0$ for $\re s>\tfrac34$.
\item \emph{A first form of the analytic half.} The identity of \cref{st:balance} and the
  reduction of \cref{sec:reduction} exist a second time, as \lean{AskMainBalance} (proved as
  \lean{colour_balance}) and \lean{GrowthInequalityForm}, with the record \lean{HostFacts}, for
  hosts ramified at most at $2$ and two odd primes. The path uses the versions that allow a set
  $S$ of primes in place of $2$ and the first odd prime (\lean{ch_balance}, \lean{ch_core},
  \lean{HostFactsS}).
\item \emph{Other definitions of the gain.} Three further definitions named \lean{gainP}, in
  other namespaces than the one of \cref{sec:adm}.
\item \emph{Other forms of single steps:} the jet statement \lean{thmBGauss}; a jet form of
  the integrality at inert primes, \lean{JetsInert}; a second form
  of the mean bound, \lean{MeanBoundSecondForm}; the determinant statements \lean{DetGain}
  (if $g_p\le v_p(\Delta)$ on a set of primes then $\sum g_p\log p\le\log\abs\Delta$),
  \lean{DetChain} and \lean{DetBlocks}; a Bessel period
  statement \lean{BesselPeriodGGP}; the congruence-number statements for $\GSp_4$
  (\lean{NumberSp4}) and the first form for $\U(3)$ (\lean{CongruenceNumberRankThree});
  and one of the derivations of the lower bound for the scalar of the lifts,
  \lean{lamLow_of_classicalTheorems}.
\end{itemize}

In the other direction, every theorem declared in the file that collects the results taken from
the literature (\cref{sec:literature}), \hubtheorems{} in all with the auxiliary ones, is in
the cone: the proof term of the final theorem refers to every one of them.

The archive declares no axiom of its own, on the path or off it: its \projnames{} declared
names are theorems, definitions, inductive types with their constructors and recursors, and
\projopaque{} opaque constants (\coneprojopaque{} of them in the cone), and nothing else.
